\documentclass[10pt]{article}
\usepackage[T1]{fontenc}
\usepackage[utf8]{inputenc}
\usepackage[french]{babel}
\usepackage{amsmath,amssymb,mathrsfs,mathpazo}
\usepackage{geometry}
\geometry{a4paper,margin=2cm}
\usepackage{graphicx,array,tabularx,tabularray,colortbl,multirow,multicol,booktabs,enumitem,pifont,ulem,fancybox,diagbox,cancel}
\usepackage{pstricks,pstricks-add,pst-plot,pst-tree,pst-node,pst-text,pst-func,pst-math,pst-bspline,pst-3dplot,pst-eucl,pst-solides3d}
\let\pscalculate\undefined
\usepackage{pst-bezier}
\usepackage{tikz,pgfplots,tkz-tab}
\usetikzlibrary{arrows,arrows.meta,patterns,calc,positioning,trees,decorations.markings,shapes.geometric,shapes.misc,shapes.arrows,intersections,decorations.pathmorphing,decorations.pathreplacing,backgrounds,shadows,babel}
\pgfplotsset{compat=1.18}
\usepgfplotslibrary{fillbetween}
\usepackage[np]{numprint}
\usepackage{listings,makecell,pas-tableur,twoopt}
\let\barre\undefined
\usepackage[pstricks]{tablvar}
\definecolor{darkgreen}{RGB}{0,100,0}
\providecommand{\red}{\color{red}}
\providecommand{\blue}{\color{blue}}
\providecommand{\green}{\color{green}}
\providecommand{\black}{\color{black}}
\providecommand{\white}{\color{white}}
\providecommand{\gray}{\color{gray}}
\providecommand{\cyan}{\color{cyan}}
\providecommand{\magenta}{\color{magenta}}
\providecommand{\yellow}{\color{yellow}}
\providecommand{\violet}{\color{violet}}
\providecommand{\orange}{\color{orange}}
\providecommand{\euro}{\texteuro}
\providecommand{\EUR}{\texteuro}
\providecommand{\ieme}{\textsuperscript{e}}
\providecommand{\degres}{\ensuremath{^\circ}}
\providecommand{\degree}{\ensuremath{^\circ}}
\providecommand{\vv}[1]{\overrightarrow{#1}}
\providecommand{\Coord}[1]{#1}
\providecommand{\decosix}{}
\providecommand{\starredbullet}{\ensuremath{\bullet}}
\pagestyle{empty}
\setlength{\parindent}{0pt}
\usepackage{adjustbox}\DeclareUnicodeCharacter{2236}{:}\DeclareUnicodeCharacter{2190}{\ensuremath{\leftarrow}}
\providecommand{\R}{\mathbb{R}}
\providecommand{\N}{\mathbb{N}}
\providecommand{\Z}{\mathbb{Z}}
\providecommand{\Q}{\mathbb{Q}}
\providecommand{\D}{\mathbb{D}}
\providecommand{\C}{\mathbb{C}}
\providecommand{\e}{{\mathrm{e}}}
\providecommand{\vect}[1]{\overrightarrow{#1}}
\providecommand{\vectt}[1]{\overrightarrow{\mathrm{#1}}}
\providecommand{\vvt}[1]{\overrightarrow{\mathrm{#1}}}
\providecommand{\np}[2][]{{#2~#1}}
\providecommand{\vv}[1]{{\overrightarrow{#1}}}
\renewcommand{\d}{{\mathrm{d}}}
\providecommand{\degres}{{^\circ}}
\providecommand{\degre}{{^\circ}}
\providecommand{\blue}{\color{blue}}
\providecommand{\red}{\color{red}}
\providecommand{\enskip}{\quad}
\providecommand{\llbracket}{\lbrack\!\lbrack}
\providecommand{\rrbracket}{\rbrack\!\rbrack}
\providecommand{\euro}{{\text{€}}}
\providecommand{\no}{{\text{n°}}}
\providecommand{\interval}[2]{{[#1;#2]}}
\providecommand{\intervFO}[2]{{[#1;#2[}}
\providecommand{\dotfill}{\ldots}

\providecommand{\ds}{\displaystyle}
\providecommand{\textding}[1]{#1}
\providecommand{\Oij}{\(\mathrm{O};\vec{\imath},\vec{\jmath}\)}
\providecommand{\Oijk}{\(\mathrm{O};\vec{\imath},\vec{\jmath},\vec{k}\)}
\providecommand{\Ouv}{\(\mathrm{O};\vec{u},\vec{v}\)}
\providecommand{\intervFF}[2]{\ensuremath{[#1;#2]}}
\renewcommand{\euro}{\ifmmode\text{\texteuro}\else\texteuro\fi}
\DeclareUnicodeCharacter{00B0}{\ifmmode^{\circ}\else\textdegree\fi}
\DeclareUnicodeCharacter{20AC}{\ifmmode\text{\texteuro}\else\texteuro\fi}
% Énoncés : Éducation nationale. Transcription et corrigés : APMEP, Denis Vergès et contributeurs. Découpage et classement : CoopMaths, Sébastien Lozano, Éric Elter et contributeurs. Adaptation de la présentation : Maths Explorer.
% CC BY-SA : https://coopmaths.fr/www/about/
% https://coopmaths.fr/alea/static/eam/2026/tex/eam_specialite_2026_06_antilles_2_cor.tex
\begin{document}
\noindent\textit{APMEP / CoopMaths. CC BY-SA. Présentation adaptée par Maths Explorer.}\par\medskip
\begin{enumerate}
	\item Comme on interroge un spectateur au hasard, on est en situation d'équiprobabilité, et les proportions données dans l'énoncé sont assimilables à des probabilités.

	On peut donc compléter l'arbre ci-dessous :

	
	\begin{center}
		\begin{tikzpicture}[yscale=-1,
			grow = right,
			level distance=3cm,
			level 1/.style={sibling distance=2cm},
			level 2/.style={sibling distance=1cm},
			edge from parent/.style={draw},
			edge from parent path=
			{(\tikzparentnode.east)--(\tikzchildnode.west)},
			every node/.style={fill=white}]
			\node {}
			child {
				node {$A$}
				child { node {$J$} edge from parent node {0,4} }
				child { node {$\overline{J}$} edge from parent node {0,6} }
				edge from parent node {0,6}
			}
			child {
				node {$\overline{A}$}
				child { node {$J$} edge from parent node {0,2} }
				child { node {$\overline{J}$} edge from parent node {0,8} }
				edge from parent node {0,4}
			};
		\end{tikzpicture}
		
		
		
			
				
				
			
			
				
				
			
			
		
	\end{center}

	\item On a, d'après la propriété des probabilités conditionnelles :

	$P(A \cap J) = P(A) \times P_{A}(J) = 0,6 \times 0,4 = 0,24$.

	Dans le contexte de l'exercice, cela signifie que 24\,\% des spectateurs sont des personnes à la fois jeunes et abonnées.

	\item Les évènements $A$ et $\overline{A}$ forment une partition de l'univers, donc, d'après la loi des probabilités totales, on a :

	$P\big(J\big) = P\big(A \cap J\big) + P\big(\overline{A} \cap J \big) = 0,24 + 0,4 \times 0,2 = 0,24 + 0,08 = 0,32$.

	On a bien trouvé $P(J) = 0,32$. Cela signifie que 32\,\% des spectateurs sont jeunes.

	\item L'affirmation peut être traduite sous forme d'inégalité mathématique :\quad $P_J(A) > \dfrac{1}{2}$.

	Or, d'après la définition des probabilités conditionnelles :\quad $P_{J}(A) = \dfrac{P(A \cap J)}{P(J)} = \dfrac{0,24}{0,32} = \dfrac{24}{32} = \dfrac{3}{4} > \dfrac{1}{2}$.

	L'affirmation est donc correcte, quand on interroge un spectateur jeune au hasard, il y a trois chances sur quatre, donc plus d'une sur deux, qu'il soit abonné.

	\item \begin{enumerate}
		\item La variable aléatoire $X$ prend trois valeurs : 5, 2 et 0.

		On a dès lors la loi de probabilité pour $X$ :

		\begin{center}
		\begin{tblr}{width=0.8\linewidth,hlines, vlines, colspec={c*3{X[c]}}}
			$k$                           & 0          & 2                       & 5              \\
			évènement $\big\{X = k\big\}$ & $A \cap J$ & $A \cap \overline{J}$   & $\overline{A}$ \\
			$P(X = k)$                    & 0,24       & $0,6 \times 0,6 = 0,36$ & 0,4
		\end{tblr}
		\end{center}

		\item On applique la formule qui définit l'espérance :

		$\mathrm{E}(X) = 0 \times P(X = 0) + 2 \times P(X = 2) + 5 \times P(X = 5) = 0 + 2 \times 0,36 + 5 \times 0,4 = 0 + 0,72 + 2 = 2,72$.

		Dans le contexte de l'exercice, cela signifie que le prix moyen payé par les spectateurs de ce cinéma est de 2,72\,\euro{}.
	\end{enumerate}
\end{enumerate}

\needspace{10\baselineskip}

\end{document}
